\subsection{Shape versus Volume: Making Large Flat Extra Dimensions Invisible --- arXiv:hep-ph/0108115}

\textbf{Authors:} Keith R. Dienes

\paragraph{主な主張}
平坦な二次元トーラスの shape modulus は Kaluza--Klein スペクトルの mass gap や level crossing を大きく変え、高次元と4次元の Planck scale の比を保ったまま KK graviton の mass gap を任意に大きな因子で増加させて large-extra-dimension 模型への実験制限を弱め得ることを示した \cite{Dienes:2001shape}。

\paragraph{新規性}
余剰次元現象論を体積 modulus だけで特徴付けるのでは不十分であり、同じ体積スケールでも複素構造・shape によって KK spectrum が大きく異なり得ることを具体的なトーラス compactification で明示した点が新しい。

\paragraph{位置付け}
large extra dimensions の実験制限を単一の「半径」ではなく volume と shape modulus を含む幾何学として再解釈する基礎論文であり、$S^1$ 型 LED neutrino/IceCube 模型を $T^2$ へ拡張して KK mass gap の制限を複素構造依存で調べる際の直接的な出発点となる。
